% !TEX program = lualatex
\documentclass[aspectratio=43,professionalfonts,handout]{beamer}
\usepackage{beamer_template}

\newcommand{\LectureNum}{04}
\newcommand{\LectureTitle}{２次関数の最大・最小}
\newcommand{\TextPages}{p.\,79-81}
\newcommand{\NextTopic}{２次関数と２次方程式（１）}
\newcommand{\NextPages}{p.\,81-83}
\newcommand{\CourseName}{基礎数学B}
\renewcommand{\TermName}{前期}

\begin{document}

% -----------------------------------------------
% スライド1: 表紙＋目標＋用語・定義
% -----------------------------------------------
\begin{frame}{\CourseName\quad 第\LectureNum 回「\LectureTitle 」}
  {\small \TermName\quad \TeacherName\quad ／\quad 教科書 \TextPages\quad
  \textbf{目標}：頂点の位置と定義域から最大値・最小値を求められる}

  \begin{block}{用語・定義}
  \begin{description}[labelwidth=5.5em]
    \item[\kw{最大値}] 関数が取り得る最も大きい値（存在しない場合もある）
    \item[\kw{最小値}] 関数が取り得る最も小さい値（存在しない場合もある）
    \item[\kw{定義域なし}] $a>0$ のとき頂点で最小、$a<0$ のとき頂点で最大
    \item[\kw{定義域あり}] 軸の位置により、端点と頂点を比較して最大・最小を決定する
  \end{description}
  \end{block}
\end{frame}

% -----------------------------------------------
% スライド2: 公式・性質
% -----------------------------------------------
\begin{frame}{公式・性質}
  \begin{block}{}
  \begin{itemize}
    \item $a>0$：頂点 $(p,\,q)$ で最小値 $q$（最大値なし）
    \item $a<0$：頂点 $(p,\,q)$ で最大値 $q$（最小値なし）
    \item 定義域 $[\alpha,\,\beta]$ あり：①軸が定義域内 → 頂点と両端を比較
    \item \phantom{定義域 $[\alpha,\,\beta]$ あり：}②軸が定義域外 → 端点のみを比較
  \end{itemize}
  \end{block}
\end{frame}

% -----------------------------------------------
% スライド3: 例1→問1
% -----------------------------------------------
\begin{frame}{例1→問1}
  \begin{exampleblock}{例題3) p.79（定義域なし）}
    （板書で解説）
  \end{exampleblock}

  \vfill

  \begin{block}{問8) p.80\quad 自分で解いてみよう}
    （ノートに解くこと）
  \end{block}
\end{frame}

% -----------------------------------------------
% スライド4: 例2→問2
% -----------------------------------------------
\begin{frame}{例2→問2}
  \begin{exampleblock}{例題4) p.80（定義域あり）}
    （板書で解説）
  \end{exampleblock}

  \vfill

  \begin{block}{問9) p.80\quad 自分で解いてみよう}
    （ノートに解くこと）
  \end{block}
\end{frame}

% -----------------------------------------------
% まとめ
% -----------------------------------------------
\begin{frame}{まとめ}
  \begin{itemize}
    \item $a>0$ のとき頂点で最小値、$a<0$ のとき頂点で最大値
    \item 定義域が限定されているときは端点と頂点を比較する
    \item 軸が定義域の内側か外側かで場合分けが必要
  \end{itemize}

  \vfill
  {\small 基本問題：150, 151, 152}
\end{frame}

\end{document}
